\documentclass[a4paper,12pt]{extarticle}
\usepackage{geometry}
\usepackage[T2A]{fontenc}
\usepackage[utf8]{inputenc}
\usepackage[main=english, russian]{babel}
\usepackage{csquotes}
\usepackage{amsmath}
\usepackage{amsthm}
\usepackage{amssymb}
\usepackage{fancyhdr}
\usepackage{setspace}
\usepackage{graphicx}
\usepackage{colortbl}
\usepackage{tikz}
\usepackage{pgf}
\usepackage{subcaption}
\usepackage{listings}
\usepackage{indentfirst}
\usepackage[
backend=biber,
style=numeric,
urldate=long,
maxbibnames=99
]{biblatex}
\addbibresource{refs.bib}
\usepackage[colorlinks,citecolor=blue,linkcolor=blue,bookmarks=false,hypertexnames=true, urlcolor=blue]{hyperref} 
\usepackage{indentfirst}
\usepackage{mathtools}
\usepackage{booktabs}
\usepackage[flushleft]{threeparttable}
\usepackage{tablefootnote}

\usepackage{chngcntr}
\counterwithin{table}{section}
\counterwithin{figure}{section}

\graphicspath{{graphics/}}

\makeatletter
\makeatother

\geometry{left=2.5cm}
\geometry{right=1.0cm}
\geometry{top=2.0cm}
\geometry{bottom=2.0cm}
\setlength{\parindent}{1.25cm}
\renewcommand{\baselinestretch}{1.5}

\newcommand{\bibref}[3]{\hyperlink{#1}{#2 (#3)}}
\addto\captionsrussian{\def\refname{Список литературы}}

\renewcommand{\theenumi}{\arabic{enumi}}
\renewcommand{\labelenumi}{\arabic{enumi}}
\renewcommand{\theenumii}{.\arabic{enumii}}
\renewcommand{\labelenumii}{\arabic{enumi}.\arabic{enumii}.}
\renewcommand{\theenumiii}{.\arabic{enumiii}}
\renewcommand{\labelenumiii}{\arabic{enumi}.\arabic{enumii}.\arabic{enumiii}.}

\begin{document}
\begin{titlepage}
	\newpage

	{\setstretch{1.0}
		\begin{center}
			NATIONAL RESEARCH UNIVERSITY\\
			HIGHER SCHOOL OF ECONOMICS
			\\
			\bigskip
			Faculty of Computer Science\\
			Bachelor’s Programme “Data Science and Business Analytics”
		\end{center}
	}

	% \vspace{2em}
	% UDC XXXXXXX % УДК нужно указывать только для исследовательсвого проекта - удалите эту строку для программного проекта
	% \vspace{4em}

	\begin{center}
		{\bf Research Project Report on the Topic:}\\
		{\bf Domain Adaptation in Computer Vision: Synthetic-to-Real Robustness Analysis}\\
	\end{center}

	\vspace{2em}

	{\bf Submitted by the Student: \vspace{2mm}}

	{\setstretch{1.1}
		\begin{tabular}{l@{\hskip 1.5cm}l}
			Indenbom Denis Evgenyevich &
			group \#\foreignlanguage{russian}{БПАД}242, 2th year of study
		\end{tabular}}

	\vspace{1em}
	{\bf Approved by the Project Supervisor: \vspace{2mm}}

	{\setstretch{1.1}
		\begin{tabular}{l@{\hskip 1.5cm}l}
			Kopylov Ivan Stanislavovich \\
			Senior Lecturer             \\
			Faculty of Computer Science, HSE University
		\end{tabular}}

	\vspace{\fill}

	\begin{center}
		Moscow 2026
	\end{center}

\end{titlepage}
\newpage
\setcounter{page}{2}

{
	\hypersetup{linkcolor=black}
	\tableofcontents
}

\newpage

\section*{Abstract}
\addcontentsline{toc}{section}{Abstract}

This project addresses the problem of domain shift in computer vision, focusing on the transition from high-quality
synthetic and studio-grade image data to real-world user-generated photographs with variable quality. A unified
training pipeline is developed to investigate domain adaptation strategies that enable effective transfer between
source and target domains. Baseline image classification models (ResNet-50 and ViT-B/16) are implemented and evaluated
for robustness under diverse environmental factors, including noise, illumination changes, shadows, blur, and viewpoint
variations.

A comparative analysis examines performance degradation across domains and assesses the effectiveness of different
adaptation approaches. The project deliverables include a reproducible training pipeline, a public repository of
experimental results, and a course work describing the methodology and findings. The outcomes contribute to a deeper
understanding of domain adaptation challenges and provide practical insights for developing robust computer vision
systems.

\section*{\foreignlanguage{russian}{Аннотация}}
\addcontentsline{toc}{section}{Аннотация}
\foreignlanguage{russian}{
	Данный проект посвящён проблеме доменного сдвига в компьютерном зрении с акцентом на переход от высококачественных синтетических и студийных изображений к пользовательским фотографиям из реального мира с вариативным качеством. В работе разрабатывается унифицированный pipeline обучения и исследуются методы доменной адаптации, обеспечивающие перенос моделей между исходным и целевым доменами. Реализуются базовые модели классификации изображений (ResNet-50 и ViT-B/16), после чего проводится их оценка устойчивости к различным внешним факторам, включая шум, изменения освещённости, тени, размытие и вариации ракурса.

	Проводится сравнительный анализ деградации качества моделей при переносе между доменами и эффективности различных стратегий адаптации. Результатом проекта являются воспроизводимый pipeline обучения, публичный репозиторий с экспериментальными результатами и курсовая работа с описанием методологии и полученных выводов. Полученные результаты способствуют более глубокому пониманию проблем доменной адаптации и формированию практических ориентиров для разработки устойчивых систем компьютерного зрения.
}

\section*{Keywords}
\addcontentsline{toc}{section}{Keywords}
Domain Adaptation, Synthetic-to-Real Transfer, Image Classification, Robustness Analysis, Computer Vision, ResNet-50, Vision Transformer, Environmental Degradation Simulation

\newpage

\section{Introduction}
\label{sec:introduction}

\subsection{Problem Relevance}
Computer vision models often fail when deployed in real-world conditions due to the domain gap between controlled
training data (synthetic/studio) and unpredictable user-generated content. This project systematically analyses the
synthetic-to-real shift, focusing on robustness against common image quality degradations.

\subsection{Problem Statement}
The research investigates the robustness of modern computer vision models under synthetic-to-real domain shift,
focusing on how environmental distortions affect classification performance. The goal is to quantify performance
degradation across domains and explore reliable ways to measure cross-domain generalization.

\subsection{Project Objectives}
\begin{enumerate}
	\item Develop a structured experimental protocol for evaluating domain shift in image classification tasks.
	\item Train and evaluate baseline models (ResNet-50 and ViT-B/16) on high-quality source domain data.
	\item Design a configurable degradation pipeline to simulate realistic domain transitions.
	\item Implement and assess domain adaptation strategies to improve target-domain generalization.
	\item Quantitatively compare baseline and adapted models across source, target, and degraded datasets.
	\item Publish a reproducible code repository and document methodology and findings in the final course project report.
\end{enumerate}

\subsection{Task Distribution}
Individual project.

\section{Literature Review}
\label{sec:litreview}

Domain adaptation methods include domain-invariant feature learning \cite{tzeng2017adversarial} and adversarial
training \cite{ganin2016domain}. Synthetic-to-real transfer often uses simulation or generative models
\cite{richter2016playing, shrivastava2017learning}. Vision Transformers \cite{dosovitskiy2020image} show strong
performance but their robustness to quality degradation is understudied. Benchmarks like ImageNet-C
\cite{hendrycks2019benchmarking} test corruption robustness but not specifically synthetic-to-real shift with
user-photography artifacts.

\section{Detailed Work Plan}
\label{sec:workplan}
The project is divided into five sequential phases.

\subsection{Phase 1: Methods Research \& Dataset Search}
\begin{itemize}
	\item Review methods for domain adaptation.
	\item Indentify suitable datasets with source domain (synthetic/studio) and target domain (user-generated).
	\item Analyze domain gaps and document key challenges.
\end{itemize}

\subsection{Phase 2: Model Implementation}
\begin{itemize}
	\item Pre-train ResNet-50 and ViT-B/16 for classification on source domain.
	\item Establish baseline metrics on validation set.
	\item Implement degradation pipelines:
	      \begin{itemize}
		      \item Gaussian noise, brightness/contrast changes.
		      \item Synthetic shadows, blur, perspective transforms.
	      \end{itemize}
\end{itemize}

\subsection{Phase 3: Appying Domain Transition methods}
\begin{itemize}
	\item Fine-tune the pre-trained ResNet-50 and ViT-B/16 with selected domain adaptation strategies.
	\item Compare at least two adaptation approaches:
	      \begin{itemize}
		      \item Feature alignment using Maximum Mean Discrepancy (MMD) or CORAL.
		      \item Domain Adversarial Training (DANN-style gradient reversal).
	      \end{itemize}
	\item Analyze convergence behavior and adaptation stability across domains.
\end{itemize}

\newpage

\subsection{Phase 4: Experiments and Analysis}
\begin{itemize}
	\item Evaluate baseline and adapted models on source and target domains.
	\item Measure overall accuracy drop between domains.
	\item Analyze feature distributions using t-SNE or PCA visualization.
	\item Compare ResNet-50 and ViT-B/16 performance under one selected degradation type.
\end{itemize}

\subsection{Phase 5: Repository and Reporting}
\begin{itemize}
	\item Create public GitHub repository with code and documentation.
	\item Write final project report.
\end{itemize}

\section{Expected Outcomes and Significance}
\label{sec:expected}

\subsection{Expected Outcomes}
\begin{enumerate}
	\item Two baseline classification models (ResNet-50 and ViT-B/16) trained on the source domain and evaluated on both source
	      and target domains.

	\item Implemented domain adaptation strategies integrated into both architectures.

	\item A configurable degradation pipeline simulating real-world image corruptions (noise, blur, illumination changes, etc.).

	\item Quantitative comparison between:
	      \begin{itemize}
		      \item Baseline models (source-only training)
		      \item Domain-adapted models
		      \item Few-shot fine-tuned models (optional scenario)
	      \end{itemize}

	\item Analysis of domain shift impact, including:
	      \begin{itemize}
		      \item Accuracy drop from source to target domain
		      \item Robustness under synthetic degradations
		      \item Feature distribution visualization (e.g., t-SNE or PCA)
	      \end{itemize}

	\item Public GitHub repository with reproducible code, trained checkpoints, and documentation.

	\item Final project report describing methodology, experimental setup, results, and key findings.
\end{enumerate}

\begin{table}[ht]
	\caption{Anticipated results format (placeholder values).}
	\label{tab:expected_results}
	\footnotesize
	\centering
	\begin{tabular}{lccc}
		\toprule
		\textbf{Test Condition}         & \textbf{ResNet-50} & \textbf{ViT-B/16} \\
		\midrule
		Source Domain (Pristine)        & 94.5\%             & 96.1\%            \\
		Target Domain (Original)        & 78.2\%             & 82.5\%            \\
		+ Gaussian Noise ($\sigma=0.1$) & 70.1\%             & 75.3\%            \\
		+ Brightness (-50\%)            & 65.4\%             & 72.8\%            \\
		+ Motion Blur                   & 58.9\%             & 68.2\%            \\
		\textbf{Combined}               & \textbf{45.2\%}    & \textbf{55.7\%}   \\
		\bottomrule
	\end{tabular}
\end{table}

\subsection{Significance}
\begin{itemize}
	\item \textbf{Scientific:} Empirical comparison of CNN and Transformer architectures under domain shift and evaluation of practical domain adaptation strategies.
	\item \textbf{Methodological:} Systematic analysis of how different adaptation mechanisms influence cross-domain generalization.
	\item \textbf{Practical:} Provides guidance on improving model robustness when deploying classifiers trained on curated datasets to real-world user-generated data.
	\item \textbf{Reproducibility:} Publicly available experimental pipeline for evaluating domain robustness and adaptation methods.
\end{itemize}

\subsection{Risks and Mitigation}
\begin{itemize}
	\item \textbf{Limited target domain data:} Use pre-trained backbones and evaluate both unsupervised adaptation and few-shot fine-tuning scenarios.
	\item \textbf{Unstable adversarial training:} Start with feature alignment methods as stable baselines and tune learning rates carefully.
	\item \textbf{Scope becoming too broad:} Restrict experiments to two architectures and a limited set of degradations.
	\item \textbf{Computational constraints:} Use smaller batch sizes, partial fine-tuning, or frozen backbones where necessary.
\end{itemize}

\newpage

\section{Methods Research \& Dataset Search}
\subsection{Methods of Domain Adaptation}
Several domain adaptation approaches have been proposed in computer vision, with Maximum Mean Discrepancy (MMD)
\cite{hu2015deep}, CORAL \cite{hu2015deep} and DANN \cite{ganin2016domain} among the most widely used in classification
tasks. This work focuses on the study and application of these strategies.

\subsubsection{MMD}
Maximum Mean Discrepancy (MMD) is a non-parametric distance metric used to measure the distribution difference between
source and target domains in a reproducing kernel Hilbert space (RKHS). The main idea is to minimize the discrepancy
between feature representations of two domains during training.

Given source domain samples $X_s=\{x_i^s\}_{i=1}^{n_s}$ and target domain samples $X_t=\{x_j^t\}_{j=1}^{n_t}$, the
squared MMD distance is defined as

\[
	\text{MMD}^2(X_s,X_t)=\left\|\frac{1}{n_s}\sum_{i=1}^{n_s}\phi(x_i^s)-\frac{1}{n_t}\sum_{j=1}^{n_t}\phi(x_j^t)\right\|_{\mathcal{H}}^2
\]

where $\phi(\cdot)$ is a feature mapping function into RKHS $\mathcal{H}$. In practice, kernel functions such as
Gaussian RBF kernel are used to compute MMD efficiently.

\subsubsection{CORAL}
Correlation Alignment (CORAL) is a domain adaptation method that aligns the second-order statistics of source and
target feature distributions by minimizing the difference between their covariance matrices.

Let $C_s$ and $C_t$ denote the covariance matrices of source and target features respectively. The CORAL loss is
defined as

\[
	\mathcal{L}_{CORAL}=\frac{1}{4d^2}\|C_s-C_t\|_F^2
\]

where $d$ is the feature dimension and $\|\cdot\|_F$ denotes the Frobenius norm. By matching covariance statistics,
CORAL reduces domain shift while keeping feature representations discriminative.

\subsubsection{DANN}
Domain-Adversarial Neural Networks (DANN) is an adversarial learning framework that aims to learn domain-invariant
feature representations using a gradient reversal layer.

The objective function consists of two parts: classification loss and domain discrimination loss:

\[
	\mathcal{L}=\mathcal{L}_y(\hat{y},y)-\lambda \mathcal{L}_d(\hat{d},d)
\]

where $\mathcal{L}_y$ is the task-specific loss (e.g., cross-entropy loss), $\mathcal{L}_d$ is the domain
classification loss, and $\lambda$ is a trade-off hyperparameter. The gradient reversal layer multiplies the gradient
by $-1$ during backpropagation to enforce domain confusion and promote domain-invariant features.

\subsection{Suitable Dataset}
\subsection{Key Challenges}

\section{Model Implementation}

% TODO: Fill this section

\newpage
\printbibliography[heading=bibintoc]

\end{document}